% !TEX root = ../main.tex
\part{Combate}

\section{Sequência de Disparo}

O combate à distância é resolvido inteiramente pelo jogador Atacante em duas etapas rápidas. O Defensor não realiza jogadas de ``esquiva'', confiando apenas em sua Cobertura e no Treinamento de suas tropas --- o mesmo princípio usado no Bolt Action.

\subsection{Passo A: Jogada de Acerto (Balística)}

O atacante reúne seu Pool de Dados e rola. O valor base para acertar é 4+. Aplique os modificadores abaixo --- todos são cumulativos entre si. Uma unidade em estado Abaixar (Down) concede -1 adicional a quem atira nela, empilhando com a cobertura do terreno que ela já tiver.

\begin{center}
\begin{tabular}{p{7cm}p{6.3cm}}
\toprule
Situação & Modificador \\
\midrule
Campo Aberto & 0 (permanece 4+) \\
Cobertura Leve (arbustos, cercas, floresta) & -1 (precisa de 5+) \\
Cobertura Pesada (muros, ruínas, dentro de casas) & -2 (precisa de 6+) \\
Bunker / Fortificação & -3 \\
Atirador se moveu (ordem Avançar) & -1 na dificuldade \\
Alvo a Longa Distância (> metade do alcance) & -1 na dificuldade \\
Alvo em estado Abaixar (Down) & -1 na dificuldade \\
Ocupação Parcial de Edifício, maioria dentro & Conta como Cobertura Pesada (-2) \\
Ocupação Parcial de Edifício, maioria fora & Conta como Cobertura Leve ou Aberto (defensor escolhe em empate) \\
\bottomrule
\end{tabular}
\end{center}

\begin{nota}
Um resultado natural de 6 no dado é sempre um acerto crítico, ignorando modificadores. Um resultado natural de 1 é sempre falha. Todos os modificadores acima são cumulativos entre si --- por exemplo, um alvo em Cobertura Leve, a Longa Distância e em Abaixar soma -3 (precisa de 7+, ou seja, só acerta com 6 natural).
\end{nota}

\subsection{Passo B: Jogada de Dano (Letalidade)}

Separe os dados que conseguiram acertar no Passo A. Role-os novamente para determinar o efeito. O valor necessário depende da qualidade da tropa alvo:

\begin{itemize}[leftmargin=1.2cm]
  \item \textbf{Alvo Inexperiente (Conscript)}: sofre baixa com 3+.
  \item \textbf{Alvo Regular (Soldado Padrão)}: sofre baixa com 4+.
  \item \textbf{Alvo Veterano (Elite/Oficiais)}: sofre baixa com 5+.
\end{itemize}

Contra veículos, o TN de dano depende do valor de Penetração (Pen) da arma versus o valor de Proteção (Prot) do alvo: role 1D6, some o Pen da arma e subtraia o Prot do veículo. Um resultado de 1 ou mais é um ferimento.

\subsection{Passo C: Resolução}

\textbf{Baixas}: para cada sucesso no Passo B, o defensor remove uma miniatura (sua escolha).

\textbf{Médico}: se houver um médico na unidade ou próximo dela, tente o salvamento agora (veja Parte IV, Médico de Combate).

\textbf{Aplicação de Stress}: se a unidade sofreu pelo menos 1 acerto de Arma Leve (mesmo sem causar baixa), ela sobe 1 nível no dado de Stress. Armas de Suporte (LMG/HMG) ou Armas Explosivas (HE) sobem 2 níveis (veja Parte III, Moral e Stress).

\section{Combate Corpo a Corpo (CQC)}

O CQC ocorre quando uma unidade executa uma ordem Avançar ou Correr bem-sucedida para entrar em contato base a base com uma unidade inimiga.

\textbf{Ataque}: ambos os lados rolam 2D6 mais um D6 para cada modelo na unidade.

\textbf{Acertos}: o número base necessário para um acerto é 4+.

\textbf{Ferimentos}: os acertos são imediatamente aplicados como ferimentos. Não há defesas no CQC.

\textbf{Resultado}: o lado que causou mais ferimentos vence o combate. O lado perdedor é destruído, e o vencedor pode consolidar 2cm.

\section{Moral e Stress (Sistema do Dado Vermelho)}

O medo de morrer é representado por um único dado d6 vermelho colocado ao lado de cada unidade sob pressão, mostrando seu Nível de Stress atual (1 a 6). Uma unidade sem stress não tem dado nenhum ao seu lado.

\subsection{Recebendo Stress}

\begin{itemize}[leftmargin=1.2cm]
  \item \textbf{Armas Leves} (Rifles, Pistolas, SMG) com pelo menos 1 acerto: sobe 1 nível.
  \item \textbf{Armas de Suporte} (LMG, HMG) ou \textbf{Armas Explosivas} (HE --- tanques, bazucas, morteiros): sobe 2 níveis.
  \item O nível nunca ultrapassa 6.
  \item \textbf{Veteranos} ignoram o 1º nível de Stress que receberiam a cada turno (a primeira vez que subiria o dado naquele turno, simplesmente não sobe).
\end{itemize}

\subsection{Efeitos por Nível}

\begin{itemize}[leftmargin=1.2cm]
  \item \textbf{Nível 1-2}: -5cm no movimento (Padrão e Corrida).
  \item \textbf{Nível 3-5}: -10cm no movimento (em vez de -5cm), -1 no Teste de Acerto, e -10cm no alcance efetivo (perde a faixa de Longo Alcance).
  \item \textbf{Nível 6}: a unidade entra em pânico, bate em retirada para fora da mesa e é removida do jogo (ver exceção da doutrina Fallschirmjäger/SS, Parte VII).
\end{itemize}

\begin{nota}
Não há mais teste de ativação para unidades com Stress --- os efeitos acima são aplicados diretamente, sem rolagem adicional.
\end{nota}

\section{Recuperação de Stress}

Existem duas maneiras de uma unidade reduzir seu Nível de Stress.

\subsection{A Ordem ``Recuperar'' (Rally)}

\textbf{Custo}: 1 PC do Sargento (ou de um Tenente próximo).

\textbf{Ação}: a unidade NÃO se move e NÃO atira neste turno.

\textbf{Resolução}: role 1D6 e reduza o Nível de Stress nesse valor (mínimo 0 --- remova o dado da mesa se chegar a 0).

\subsection{Intervenção do Oficial}

O Tenente pode gastar 1 PC para ordenar ``Recuperar'' em uma unidade a até 15cm, seguindo as mesmas regras acima. Útil quando o Sargento da unidade morreu.
